Алгоритмы сортировки "--- фундаментальное понятие в компьютерных науках. Они широко применяются в системах управления базами данных, поисковых системах, обработке больших массивов информации и многих других областях. Эффективность сортировки напрямую влияет на производительность программного обеспечения, особенно при работе с большими объёмами данных.

С развитием языков программирования и компиляторов возникает вопрос: насколько существенно различие в производительности реализации одних и тех же алгоритмов на разных языках? Для ответа на этот вопрос в данной работе выбраны два популярных языка:  Python и C++. Python известен своей простотой и богатой экосистемой, но часто уступает в скорости выполнения. C++ предоставляет больше контроля над памятью и позволяет писать очень быстрый код, но требует более аккуратного подхода.

\textbf{Цель работы} "--- провести сравнительный анализ трёх классических алгоритмов сортировки (пузырьковой, быстрой и сортировки слиянием), реализованных на Python и C++, оценить их временную сложность на практике и сопоставить с теоретическими оценками.

\textbf{Задачи исследования:}
\begin{enumerate}
    \item изучить теоретические основы выбранных алгоритмов, включая формулы временной сложности;
    \item выполнить реализацию каждого алгоритма на обоих языках программирования, следуя принципам структурного программирования;
    \item разработать методику тестирования, включая генерацию случайных массивов различного размера ($n = 1000, 5000, 10000, 50000$);
    \item провести серию замеров времени выполнения (не менее 5 повторений для каждого размера) и вычислить средние значения;
    \item визуализировать полученные результаты в виде графиков и таблиц;
    \item проанализировать различия в производительности и объяснить их с точки зрения особенностей компиляции и интерпретации.
\end{enumerate}

\textbf{Объект исследования} "--- сортировка массивов целых чисел. 
\textbf{Предмет исследования} "--- скорость выполнения в зависимости от языка и размера данных.

\textbf{Практическая значимость} "--- результаты помогут разработчикам выбирать язык для ресурсоёмких задач. \cite{knuth1998, cormen2009, python_sort}
